\chapter{Mengukur dan Waktu}\label{ch:g1:measure}

Pensil mana yang lebih panjang? Hari ini hari apa? Berapa harga satu
permen? Bab ini membandingkan panjang dan mengukurnya dalam
sentimeter, mengurutkan hari dalam seminggu, membaca jam bulat pada
jam dinding, dan membilang uang logam kecil.

\section{Lebih panjang dan lebih pendek}

\begin{method}[Membandingkan panjang]\label{met:g1:measure:compare}
Untuk membandingkan dua panjang, sejajarkan keduanya dengan
\emph{satu ujung yang berimpit}:
\begin{enumerate}
  \item letakkan kedua benda berdampingan, mulai dari garis yang sama;
  \item benda yang menjulur lebih jauh itulah yang \emph{lebih
        panjang}, yang lain \emph{lebih pendek};
  \item jika keduanya sampai di tempat yang sama, keduanya
        \emph{sama panjang}.
\end{enumerate}
Menyamakan titik awal itulah seluruh rahasianya --- membandingkan
dari awal yang tidak rata akan menipu mata.
\end{method}

\begin{omfigure}
\begin{tikzpicture}[scale=0.9]
  \draw[very thick, omDef] (0,0.6) -- (3.4,0.6);
  \draw[very thick, omThm] (0,0) -- (2.2,0);
  \draw[gray, dashed] (0,-0.2) -- (0,0.8);
  \node[right, font=\small] at (3.6,0.6) {lebih panjang};
  \node[right, font=\small] at (3.6,0) {lebih pendek};
\end{tikzpicture}

{\small Dua lidi yang diimpitkan di sebelah kiri: yang biru menjulur
lebih jauh, jadi ia lebih panjang.}
\end{omfigure}

\begin{definition}[Sentimeter]\label{def:g1:measure:cm}
Untuk mengatakan \emph{berapa} panjangnya, kita mengukur dengan
penggaris dalam \emph{sentimeter} (cm). Letakkan angka $0$ penggaris
di awal benda, lalu bacalah bilangan di ujungnya.
\end{definition}

\begin{omfigure}
\begin{tikzpicture}[scale=1.0]
  % badan penggaris menjulur sedikit melewati tanda terakhir agar tidak ada angka yang terpotong
  \draw[thick] (-0.4,0) rectangle (8.4,0.7);
  \foreach \i in {0,...,8} {
    \draw[thick] (\i,0.7) -- (\i,0.4);
    \node[font=\footnotesize] at (\i,0.2) {\i};
  }
  \draw[very thick, omThm] (0,1.15) -- (5,1.15);
  \node[omThm, above, font=\small] at (2.5,1.22) {pensilnya: $5$ cm};
\end{tikzpicture}

{\small Mengukur: angka $0$ di titik awal, lalu bacalah ujungnya.
Pensil ini panjangnya $5$ cm. (Milimeter dan meter menyusul pada
\cref{ch:g3:measure}.)}
\end{omfigure}

\section{Hari dalam seminggu}

\begin{example}[Tujuh hari]\label{ex:g1:measure:days}
Satu minggu punya $7$ hari, selalu dengan urutan yang sama:
\[
\text{Senin, Selasa, Rabu, Kamis, Jumat, Sabtu,
Minggu.}
\]
Sesudah Minggu, Senin mulai lagi. \emph{Kemarin} adalah hari
sebelumnya, \emph{besok} hari sesudahnya: jika hari ini Rabu, kemarin
hari Selasa dan besok hari Kamis.
\end{example}

\section{Jam bulat pada jam dinding}

\begin{method}[Membaca jam bulat]\label{met:g1:measure:clock}
Ketika jarum panjang menunjuk lurus ke atas ke angka $12$, saat itu
jam bulat: bacalah angka yang ditunjuk jarum \emph{pendek}. Jarum
pendek di $3$, jarum panjang di $12$: saat itu pukul $3$.
\end{method}

\begin{omfigure}
\begin{tikzpicture}[scale=0.95]
  \draw[thick] (0,0) circle (1.3);
  \foreach \h in {1,...,12} {
    \node[font=\footnotesize] at ({90-\h*30}:1.05) {\h};
    \draw ({90-\h*30}:1.22) -- ({90-\h*30}:1.3);
  }
  \draw[very thick, omDef] (0,0) -- (0:0.62);
  \draw[thick, omThm] (0,0) -- (90:1.02);
  \node[below, font=\small] at (0,-1.55) {pukul $3$};
  \begin{scope}[xshift=4.4cm]
  \draw[thick] (0,0) circle (1.3);
  \foreach \h in {1,...,12} {
    \node[font=\footnotesize] at ({90-\h*30}:1.05) {\h};
    \draw ({90-\h*30}:1.22) -- ({90-\h*30}:1.3);
  }
  \draw[very thick, omDef] (0,0) -- ({90-8*30}:0.62);
  \draw[thick, omThm] (0,0) -- (90:1.02);
  \node[below, font=\small] at (0,-1.55) {pukul $8$};
  \end{scope}
\end{tikzpicture}

{\small Jam bulat: jarum panjang merah di angka $12$, jarum pendek
biru menunjukkan jamnya. (Setengah jam dan seperempat jam menyusul
pada \cref{ch:g2:measure}.)}
\end{omfigure}

\section{Uang logam kecil}

\begin{example}[Menyusun sejumlah uang]\label{ex:g1:measure:coins}
Dengan uang logam $1$, $2$ dan $5$, banyak \omterm{def:g1:addition:def}{jumlah} dapat disusun
dengan beberapa cara. Untuk menyusun $7$:
\[
5 + 2, \qquad
5 + 1 + 1, \qquad
2 + 2 + 2 + 1 .
\]
Memakai logam terbesar lebih dulu biasanya cara tercepat, dan
memakai logam paling sedikit.
\end{example}

\begin{omfigure}
\begin{tikzpicture}[scale=0.85]
  \draw[thick, fill=omMeth!20] (0,0) circle (0.5);
  \node[font=\small] at (0,0) {$5$};
  \node[font=\large] at (1.0,0) {$+$};
  \draw[thick, fill=omDef!20] (2.0,0) circle (0.42);
  \node[font=\small] at (2.0,0) {$2$};
  \node[font=\large] at (3.0,0) {$=$};
  \node[font=\small] at (4.0,0) {$7$};
\end{tikzpicture}

{\small Satu logam lima dan satu logam dua menjadi tujuh.}
\end{omfigure}

\section{Latihan}

\begin{exercise}[$\star$]\label{exo:g1:measure:1}
Sejajarkan tiga pensil dengan satu ujung berimpit. Mana yang paling
panjang? Mana yang paling pendek?
\end{exercise}

\begin{exercise}[$\star$]\label{exo:g1:measure:2}
Ukurlah dengan penggarismu: sisi panjang buku ini; jari kelingkingmu;
sebuah sendok. Jawablah dalam sentimeter bulat.
\end{exercise}

\begin{exercise}[$\star$]\label{exo:g1:measure:3}
Sebutkan harinya: sesudah Senin; sebelum Minggu; dua hari sesudah Jumat.
\end{exercise}

\begin{exercise}[$\star$]\label{exo:g1:measure:4}
Hari ini hari Kamis. Kemarin hari apa? Besok hari apa? Tiga hari lagi
hari apa?
\end{exercise}

\begin{exercise}[$\star$]\label{exo:g1:measure:5}
Gambarlah jam yang menunjukkan pukul $6$. Di mana jarum pendeknya? Di
mana jarum panjangnya?
\end{exercise}

\begin{exercise}[$\star$]\label{exo:g1:measure:6}
Bacalah jam bulatnya: jarum pendek di $5$, jarum panjang di $12$;
jarum pendek di $11$, jarum panjang di $12$.
\end{exercise}

\begin{exercise}[$\star$]\label{exo:g1:measure:7}
Susunlah $8$ dengan uang logam $1$, $2$ dan $5$, dengan dua cara berbeda.
\end{exercise}

\begin{exercise}[$\star$]\label{exo:g1:measure:8}
Rudi punya satu logam $5$ dan tiga logam $2$. Berapa \omterm{def:g1:addition:def}{jumlah} uangnya
seluruhnya?
\end{exercise}

\begin{exercise}[$\star\star$]\label{exo:g1:measure:9}
Sebuah permen berharga $3$. Zahra membayar dengan logam $5$. Berapa
kembaliannya? (Bilanglah maju dari $3$ ke $5$.)

\end{exercise}
